
\subsubsection{Main Results}

Table 1 presents gate metrics against success criteria. All three metrics failed.

\textbf{Table 1: Gate Metric Results}

\begin{table}[htbp]
\centering
\begin{tabular}{lllll}
\toprule
Metric & Target & Actual & Status & Gap \\
\midrule
Reconstruction Error & <10.0\% & 19.18\% & FAIL & +9.18pp \\
Frozen-K Generalization & <10.0\% & 10.31\% & FAIL & +0.31pp \\
Kernel Robustness & ≥90.0\% & 0.00\% & FAIL & -90.0pp \\
\bottomrule
\end{tabular}
\end{table}

\textbf{Key Observations}:

\begin{enumerate}
\item \textbf{Kernel robustness complete failure (0.00\%)}: The most critical finding. Despite explicit equivariance loss training with λ_equiv=0.5, the encoder achieved 0.00\% permutation invariance (target ≥90\%). Random weight permutations cause large divergences in representations: D = ||E(g · w) - E(w)|| > 0.01 for 100\% of tested permutations. This demonstrates that MSE-based equivariance loss L_equiv = ||E(g · M) - ρ(g)E(M)||² in the tested configuration does not enforce group homomorphism structure—gradient descent on reconstruction objectives does not learn permutation group structure in this setting.
\end{enumerate}

\begin{enumerate}
\item \textbf{Reconstruction error indicates insufficient dimensionality (19.18\%)}: The K=32 space cannot capture structure from 1000-dimensional weight vectors, showing 19.18\% information loss (target <10\%, gap +9.18pp). This suggests dimensionality requirements are underestimated. Extrapolating to real 100K-dimensional models, this pattern suggests substantially higher K may be necessary, though the exact relationship remains unknown.
\end{enumerate}

\begin{enumerate}
\item \textbf{Frozen-K generalization marginally fails (10.31\%)}: The encoder trained on CNN+Transformer shows 10.31\% reconstruction error on held-out RNN models (target <10\%, gap +0.31pp). While the 0.31pp gap is small, when combined with visualization evidence showing architecture-specific clustering (Figure 2), this suggests the encoder may learn family-specific rather than shared representations. The architecture embeddings designed to provide helpful context may anchor representations to family-specific coordinates.
\end{enumerate}

\subsubsection{Training Dynamics}

Figure 1 shows training curves over 12 epochs (early stopped from planned 20 epochs). The combined loss plateaus at epoch 8, triggering early stopping at epoch 12 with patience=10.

\textbf{Analysis}: Early stopping at epoch 12 occurred due to validation loss plateau. The reconstruction loss improves steadily (0.25 → 0.19) while equivariance loss remains high and volatile (0.15-0.20 range). This suggests the two loss objectives may have different optimization dynamics, which could indicate conflicting gradients, though further investigation would be needed to confirm this.

\subsubsection{Representation Structure}

Figure 2 visualizes the learned space via t-SNE projection, coloring points by architecture family (blue=CNN, orange=Transformer, green=RNN). The space shows strong architecture-specific clustering rather than shared structure. CNNs, Transformers, and RNNs form distinct clusters with minimal overlap. This confirms the frozen-K generalization failure: the encoder learns three separate coordinate systems rather than factoring out architecture conventions into a unified space.

The architecture embeddings (64-dim) designed to help cross-architecture learning may instead cause this clustering. By providing architecture family context, we may signal the encoder to learn family-specific representations. A pure architecture-agnostic encoder without family embeddings should be tested.

\subsubsection{Error Distribution}

Figure 3 shows the distribution of reconstruction errors across the 150-model test set. Reconstruction errors range from 12\% (best case) to 35\% (worst case) with mean 19.18\% ± 4.3\%. No model achieves the <10\% target. The right-skewed distribution with long tail suggests certain model types are particularly poorly represented in the K=32 space.

Even best-case performance (12\%) fails to meet the target (<10\%), confirming that K=32 is insufficient rather than requiring better optimization. The high-variance tail (some models >25\% error) suggests the space cannot uniformly represent the diversity of architectures.
